% Clase 20 — Belmont Muñoz Samuel (botones y enlace USB HID)
\paragraph{Botones y enlace USB HID}
\textbf{Autor:} Belmont Muñoz Samuel

La imagen de los controles es \texttt{controls\_report\_t} (\texttt{controls\_report.h}, líneas 17 a 20):
cuatro ejes y los botones, todos \texttt{uint16\_t}. Los 10 botones de los dos 74HC165 no caben en 8 bits,
por eso \texttt{buttons} es de 16; los 6 bits de sobra se limpian con la máscara \texttt{0x03FF} antes de enviar
(\texttt{usb\_data\_tx.c}, línea 35). La estructura mide $4 \times 2 + 2 = 10$ bytes sin relleno, y el compilador
lo verifica con \texttt{\_Static\_assert(sizeof(controls\_report\_t) == 10)} (\texttt{controls\_report.h}, línea 31).
Esos 10 bytes son exactamente el reporte de entrada HID que recibe la Pi, así que no hay que convertir nada antes de enviar.

\texttt{volatile} está en \texttt{g\_controls} (\texttt{main.c}, línea 19) porque la escriben los DMA de la PIO y del
ADC sin pasar por la CPU: el compilador no ve ninguna escritura en el código y, sin \texttt{volatile}, podría
quedarse con un valor viejo en un registro. Por lo mismo, \texttt{controls\_snapshot()} lee campo por campo con un
puntero \texttt{volatile uint16\_t} (\texttt{usb\_data\_tx.c}, línea 29), y el búfer del \textit{endpoint} en la DPSRAM
también es \texttt{volatile} (línea 225), porque lo lee el controlador USB por hardware.

\textbf{Tamaño del firmware} (Pico SDK 2.1.1, \texttt{arm-none-eabi-size}):
\begin{lstlisting}
   text    data     bss     dec     hex filename
  23468       0    3520   26988    696c arcade_ctrl.elf
\end{lstlisting}
Ocupa 23\,468 B de Flash (1.1\,\% de 2 MB) y 3\,520 B de RAM estática (1.3\,\% de 264 KB).
